% SPDX-License-Identifier: CC-BY-4.0
% Copyright 2025 Florian Aram Feuerriegel — kassensturz.org
% Abschnitt 4 — Wirkungen der Stellgrößen. Jede Zahl im Text ist ein Makro aus
% generiert/werte.tex oder generiert/makros.tex (\W, \Wb), jede Tabelle erzeugt
% messung.js.

\section{Wirkungen der Stellgrößen}
\label{sec:wirkungen}

\subsection{Messung}
\label{sec:messung}

Die Wirkung einer Stellgröße ist gemessen, nicht hergeleitet. Ausgehend vom
Status quo wird sie um einen festen Schritt verschoben
(Tabelle~\ref{tab:stellgroessen}), das Modell rechnet den Pfad über alle fünf
Perioden neu, und die Differenz zum Status-quo-Pfad ist die Wirkung. Sätze
verschieben sich um einen Prozentpunkt, Vermögen- und Bodenwertsteuer um
0,1~Punkte, monatliche Beträge um 10\,\euro, der Grundfreibetrag um
1.000\,\euro, die Beitragsbemessungsgrenze um 10.000\,\euro{} und der
\COz-Preis um 10\,\euro/t; die Grenze der Spitzenzone sinkt um 10.000\,\euro,
das Klimageld wird eingeschaltet. Das Modell ist stetig, aber nicht
überall differenzierbar: Knicke liegen an den Tarifgrenzen, an den
Klammern der Verhaltensreaktionen, dort, wo eine Beitragsbemessungsgrenze das
Arbeitseinkommen eines Haushaltstyps kreuzt, und an der Anrechnungsschwelle der
Mindeststeuer. Die Messwerte sind deshalb endliche Differenzen für genau diesen
Schritt, keine Ableitungen; Abschnitt~\ref{sec:nichtlinear} zeigt, wie weit sie
tragen.

Zwei Arten von Gewichten verbinden Stellgrößen und Kennzahlen. Von der
Stellgröße zu den Kanälen, also zu Aufkommen je Steuerart, Transfers, Renten,
Arbeitsangebot, Investitionen und Nachfrage, sind sie gemessen. Von den Kanälen
zu den Kennzahlen folgen sie aus den Gleichungen des Abschnitts~\ref{sec:modell}
im Status quo: Ein Euro Aufkommen verbessert den Saldo um einen Euro, ein Euro
Nachfrage bringt \Wstabil\,\euro{} Einnahmen (Abbildung~\ref{fig:struktur}).
Innerhalb einer Periode ist dieser Graph kreisfrei; Kreise entstehen erst über
die Periodengrenze, wenn Schulden Zinsen erzeugen und das BIP-Niveau die
Einnahmen der nächsten Periode skaliert.

\begin{table}[tbp]
\centering\small
\caption{Stellgrößen, Status quo und Messschritt. Die oberen stehen am
Verhandlungstisch des Planspiels, geordnet nach Ressort; die weiteren sind
keinem Ressort zugeordnet. Körperschaftsteuer und Rentenbeitrag folgen
zusätzlich dem geltenden Recht (Abschnitt~\ref{sec:aufbau}).}
\label{tab:stellgroessen}
% Erzeugt von docs/modell/messung.js — nicht von Hand bearbeiten
\begin{tabular}{@{}llrrr@{}}
\toprule
Stellgröße & Ressort & Status quo & Spielraum & Schritt \\
\midrule
Grundfreibetrag & Finanzen & 12.348\,\euro{} & 0--30.000\,\euro{} & +1.000\,\euro{} \\
Eingangssteuersatz & Finanzen & 14\,\% & 0--40\,\% & +1\,Pp. \\
Spitzensteuersatz & Finanzen & 45\,\% & 20--75\,\% & +1\,Pp. \\
Umsatzsteuer & Finanzen & 19\,\% & 0--30\,\% & +1\,Pp. \\
Erbschaftsteuer & Finanzen & 20\,\% & 0--60\,\% & +1\,Pp. \\
Körperschaftsteuer & Wirtschaft & 15\,\% & 0--40\,\% & +1\,Pp. \\
Gewerbesteuer & Wirtschaft & 14\,\% & 0--30\,\% & +1\,Pp. \\
RV-Beitrag & Soziales & 18,6\,\% & 10--30\,\% & +1\,Pp. \\
KV-Beitrag & Soziales & 17,5\,\% & 10--25\,\% & +1\,Pp. \\
Rentenniveau & Soziales & 48\,\% & 40--53\,\% & +1\,Pp. \\
Bürgergeld & Soziales & 563\,\euro{}/Monat & 0--1.500\,\euro{}/Monat & +10\,\euro{}/Monat \\
Kindergeld & Soziales & 259\,\euro{}/Monat & 0--1.000\,\euro{}/Monat & +10\,\euro{}/Monat \\
CO\textsubscript{2}-Preis & Umwelt & 55\,\euro{}/t & 0--300\,\euro{}/t & +10\,\euro{}/t \\
Klimageld & Umwelt & nein & ja/nein & ein \\
\addlinespace
\multicolumn{5}{@{}l}{\emph{Weitere Stellgrößen}} \\
Grenze Spitzensatz & -- & 277.826\,\euro{} & -- & $-$10.000\,\euro{} \\
Umsatzsteuer, ermäßigt & -- & 7\,\% & -- & +1\,Pp. \\
Abgeltungsteuer & -- & 25\,\% & -- & +1\,Pp. \\
Beitragsbemessungsgrenze & -- & 101.400\,\euro{} & -- & +10.000\,\euro{} \\
Vermögensteuer & -- & 0\,\% & -- & +0,1\,Pp. \\
Bodenwertsteuer & -- & 0,4\,\% & -- & +0,1\,Pp. \\
Mindeststeuer Milliardäre & -- & 0\,\% & -- & +1\,Pp. \\
\bottomrule
\end{tabular}

\end{table}

\subsection{Zielkonflikte}
\label{sec:zielkonflikte}

Tabelle~\ref{tab:gesamt} fasst die Wirkungen der Stellgrößen am Tisch zusammen,
in der ersten Periode und am Ende des Pfads. Die Färbung zeigt, ob sich eine
Kennzahl in die erwünschte Richtung bewegt: höherer Saldo, niedrigerer Gini,
niedrigere Armut und Emissionen, höheres BIP und Nettoeinkommen, niedrigere
Schuldenquote. Keine Zeile ist durchgehend blau. Der Gini-Koeffizient steht hier
und im Folgenden in Punkten, ein Punkt ist 0,01. Tabelle~\ref{tab:preise} bezieht
die Wirkungen auf eine Milliarde Saldoverbesserung: Jede Stellgröße wird in der
Richtung verschoben, die den Saldo verbessert, und die Änderungen von Gini und
BIP werden durch die Saldoverbesserung geteilt. Das Ergebnis ist der Preis
einer Milliarde Konsolidierung in Verteilung und Wachstum. Drei Muster treten
hervor.

\emph{Einnahmen und Gleichheit stehen im Konflikt.} Beitragssätze,
Umsatzsteuer und \COz-Preis erhöhen den Gini-Koeffizienten, ebenso jede
Kürzung von Renten, Bürger- oder Kindergeld. Steuern auf hohe Einkommen,
Kapitaleinkommen und Vermögen senken ihn, ebenso eine höhere
Beitragsbemessungsgrenze und sogar ein höherer Eingangssteuersatz, weil D1
und D2 unter dem Grundfreibetrag liegen. Die ergiebigsten Instrumente liegen
auf der regressiven Seite: Ein Prozentpunkt Rentenbeitrag verbessert den Saldo
um \Wb{kennzahl}{rv}{saldo}~Mrd.~\euro, ein Prozentpunkt Spitzensteuersatz um
\Wb{kennzahl}{spitze}{saldo}~Mrd.~\euro.

\emph{Konsolidierung kostet Wachstum, je nach Träger verschieden viel.} Weil
die Grenzkonsumneigung unten höher ist, senkt eine Milliarde, die untere und
mittlere Einkommen tragen, die Nachfrage stärker als eine Milliarde aus dem
obersten Prozent. In der ersten Periode kostet eine Milliarde über den
Rentenbeitrag \Wb{preis}{rv}{bip}~Mrd.~\euro{} BIP, über die Mindeststeuer
auf Milliardärsvermögen \Wb{preis}{zucman}{bip}~Mrd.~\euro. Was davon bleibt,
entscheidet die Angebotsseite: Das BIP-Niveau 2041 bewegen nur Stellgrößen, die
über das Arbeitsangebot oder das private Kapital wirken. Ein Prozentpunkt
Körperschaftsteuer senkt es um \Wb{pfad}{kst}{bip}~Mrd.~\euro{} und verbessert
den Saldo der letzten Periode nur um \Wb{pfad}{kst}{saldo}~Mrd.~\euro.

\emph{Transfers wirken zielgenau, aber teuer.} Ein Klimageld zahlt jedem
Haushalt \Wb{dezil}{klimageld}{D1}\,\euro{} im Jahr, senkt den Gini um
\Wb{kennzahl}{klimageld}{gini}~Punkte und kostet
\Wb{kennzahl}{klimageld}{saldo}~Mrd.~\euro. Ein Prozentpunkt Rentenniveau
senkt den Gini um \Wb{kennzahl}{rentenniveau}{gini}~Punkte, weil die gesetzliche
Rente in der unteren Hälfte einen großen Teil des Einkommens ausmacht. Je
Milliarde wirkt das Bürgergeld am stärksten: Je 10~Mrd.~\euro{} Einsparung steigt
der Gini um \Wb{preis}{bg}{gini10}~Punkte, beim Rentenniveau um
\Wb{preis}{rentenniveau}{gini10}. Ein höherer Grundfreibetrag senkt zwar den Gini,
erhöht aber die Armutsrisikoquote um \Wb{kennzahl}{freibetrag}{armut}~Prozentpunkte:
D1 und D2 zahlen keine Einkommensteuer und gewinnen nichts, während mit dem
Median die Armutsschwelle steigt.

\begin{table}[tbp]
\centering\scriptsize
\setlength{\tabcolsep}{2.8pt}
\caption{Gesamtmatrix der Wirkungen je Schritt: erste Periode (Jahreswerte
2025--2028) und Ende des Pfads, Schuldenquote und BIP-Niveau zu Beginn der
letzten Periode (2041), Emissionen 2025--2040 kumuliert; die Stellgröße bleibt
über alle Perioden verschoben.
\textcolor{pos}{Blau}: erwünschte Richtung, \textcolor{neg}{orange}: unerwünschte,
Sättigung nach Betrag relativ zum größten Wert der Spalte; grau: keine messbare
Wirkung. Gini in Punkten (1~Punkt = 0,01), Nettoeinkommen je Haushalt und Jahr,
BIP-Niveau zu Beginn der letzten Periode ohne deren Nachfrage.}
\label{tab:gesamt}
% Erzeugt von docs/modell/messung.js — nicht von Hand bearbeiten
\begin{tabular}{@{}l*{11}{r}@{}}
\toprule
 & \multicolumn{8}{c}{Erste Periode, 2025--2028} & \multicolumn{3}{c@{}}{Ende des Pfads} \\
\cmidrule(lr){2-9}\cmidrule(l){10-12}
Stellgröße (Schritt) & Saldo & Gini & Armut & Emiss. & BIP & D1 & D5 & D10c & Schulden & BIP-Niv. & CO\textsubscript{2} kum. \\
 & Mrd.\,\euro{} & Pkt. & Pp. & Mt & Mrd.\,\euro{} & \euro{} & \euro{} & \euro{} & Pp. & Mrd.\,\euro{} & Mt \\
\midrule
Grundfreibetrag (+1.000\,\euro{}) & \cellcolor{neg!45}$-$6,9 & \cellcolor{pos!35}$-$0,12 & \cellcolor{neg!70}0,29 & \textcolor{grau}{0} & \cellcolor{pos!60}8,5 & \textcolor{grau}{0} & \cellcolor{pos!69}357 & \cellcolor{pos!24}637 & \cellcolor{neg!25}0,87 & \cellcolor{pos!70}15,3 & \textcolor{grau}{0} \\
Eingangssteuersatz (+1\,Pp.) & \cellcolor{pos!25}2,6 & \cellcolor{pos!19}$-$0,04 & \cellcolor{pos!29}$-$0,09 & \textcolor{grau}{0} & \cellcolor{neg!28}$-$2,8 & \textcolor{grau}{0} & \cellcolor{neg!28}$-$100 & \cellcolor{neg!18}$-$322 & \cellcolor{pos!20}$-$0,52 & \cellcolor{neg!26}$-$3,7 & \textcolor{grau}{0} \\
Spitzensteuersatz (+1\,Pp.) & \cellcolor{pos!13}0,1 & \cellcolor{pos!16}$-$0,02 & \textcolor{grau}{0} & \textcolor{grau}{0} & \cellcolor{neg!13}$-$0,1 & \textcolor{grau}{0} & \textcolor{grau}{0} & \cellcolor{neg!42}$-$1.609 & \cellcolor{pos!13}$-$0,04 & \cellcolor{neg!12}$-$0,1 & \textcolor{grau}{0} \\
Umsatzsteuer (+1\,Pp.) & \cellcolor{pos!49}7,7 & \cellcolor{neg!22}0,06 & \cellcolor{pos!19}$-$0,03 & \textcolor{grau}{0} & \cellcolor{neg!43}$-$5,5 & \cellcolor{neg!29}$-$87 & \cellcolor{neg!44}$-$204 & \cellcolor{neg!41}$-$1.555 & \cellcolor{pos!48}$-$2,44 & \textcolor{grau}{0} & \textcolor{grau}{0} \\
Erbschaftsteuer (+1\,Pp.) & \cellcolor{pos!13}0,3 & \cellcolor{pos!13}$-$0,01 & \textcolor{grau}{0} & \textcolor{grau}{0} & \cellcolor{neg!13}$-$0,1 & \textcolor{grau}{0} & \cellcolor{neg!12}$-$2 & \cellcolor{neg!15}$-$177 & \cellcolor{pos!13}$-$0,08 & \textcolor{grau}{0} & \textcolor{grau}{0} \\
Körperschaftsteuer (+1\,Pp.) & \cellcolor{pos!22}2,0 & \cellcolor{pos!18}$-$0,04 & \cellcolor{pos!14}$-$0,01 & \textcolor{grau}{0} & \cellcolor{neg!18}$-$1,0 & \cellcolor{neg!14}$-$8 & \cellcolor{neg!16}$-$27 & \cellcolor{neg!42}$-$1.616 & \cellcolor{pos!17}$-$0,33 & \cellcolor{neg!36}$-$6,5 & \textcolor{grau}{0} \\
Gewerbesteuer (+1\,Pp.) & \cellcolor{pos!31}3,9 & \cellcolor{pos!25}$-$0,07 & \cellcolor{pos!16}$-$0,02 & \textcolor{grau}{0} & \cellcolor{neg!23}$-$2,0 & \cellcolor{neg!15}$-$16 & \cellcolor{neg!20}$-$52 & \cellcolor{neg!70}$-$3.128 & \cellcolor{pos!26}$-$0,93 & \cellcolor{neg!36}$-$6,5 & \textcolor{grau}{0} \\
RV-Beitrag (+1\,Pp.) & \cellcolor{pos!68}11,7 & \cellcolor{neg!24}0,06 & \cellcolor{pos!33}$-$0,10 & \textcolor{grau}{0} & \cellcolor{neg!70}$-$10,3 & \cellcolor{neg!36}$-$129 & \cellcolor{neg!70}$-$366 & \cellcolor{neg!30}$-$987 & \cellcolor{pos!68}$-$3,76 & \textcolor{grau}{0} & \textcolor{grau}{0} \\
KV-Beitrag (+1\,Pp.) & \cellcolor{pos!70}12,1 & \cellcolor{neg!34}0,12 & \cellcolor{pos!33}$-$0,10 & \textcolor{grau}{0} & \cellcolor{neg!66}$-$9,6 & \cellcolor{neg!36}$-$129 & \cellcolor{neg!70}$-$366 & \cellcolor{neg!25}$-$679 & \cellcolor{pos!70}$-$3,90 & \textcolor{grau}{0} & \textcolor{grau}{0} \\
Rentenniveau (+1\,Pp.) & \cellcolor{neg!40}$-$5,8 & \cellcolor{pos!51}$-$0,22 & \cellcolor{pos!13}$-$0,01 & \textcolor{grau}{0} & \cellcolor{pos!39}4,8 & \cellcolor{pos!35}122 & \cellcolor{pos!50}238 & \cellcolor{pos!14}117 & \cellcolor{neg!42}2,01 & \textcolor{grau}{0} & \textcolor{grau}{0} \\
Bürgergeld (+10\,\euro{}/Monat) & \cellcolor{neg!14}$-$0,4 & \cellcolor{pos!20}$-$0,04 & \cellcolor{pos!27}$-$0,08 & \textcolor{grau}{0} & \cellcolor{pos!14}0,4 & \cellcolor{pos!26}72 & \textcolor{grau}{0} & \textcolor{grau}{0} & \cellcolor{neg!14}0,11 & \textcolor{grau}{0} & \textcolor{grau}{0} \\
Kindergeld (+10\,\euro{}/Monat) & \cellcolor{neg!20}$-$1,7 & \cellcolor{pos!24}$-$0,07 & \cellcolor{pos!16}$-$0,02 & \textcolor{grau}{0} & \cellcolor{pos!19}1,3 & \cellcolor{pos!20}45 & \cellcolor{pos!21}59 & \cellcolor{pos!12}11 & \cellcolor{neg!20}0,52 & \textcolor{grau}{0} & \textcolor{grau}{0} \\
CO\textsubscript{2}-Preis (+10\,\euro{}/t) & \cellcolor{pos!22}2,0 & \cellcolor{neg!16}0,02 & \cellcolor{pos!14}$-$0,01 & \cellcolor{pos!70}$-$9,8 & \cellcolor{neg!20}$-$1,4 & \cellcolor{neg!17}$-$24 & \cellcolor{neg!21}$-$55 & \cellcolor{neg!18}$-$303 & \cellcolor{pos!19}$-$0,46 & \textcolor{grau}{0} & \cellcolor{pos!70}$-$103 \\
Klimageld (ein) & \cellcolor{neg!58}$-$9,7 & \cellcolor{pos!70}$-$0,32 & \cellcolor{pos!43}$-$0,15 & \textcolor{grau}{0} & \cellcolor{pos!56}7,8 & \cellcolor{pos!70}307 & \cellcolor{pos!61}307 & \cellcolor{pos!18}307 & \cellcolor{neg!45}2,19 & \textcolor{grau}{0} & \textcolor{grau}{0} \\
\bottomrule
\end{tabular}

\end{table}

\begin{table}[tbp]
\centering\footnotesize
\caption{Preis einer Milliarde Konsolidierung. Jede Stellgröße in der Richtung,
die den Saldo verbessert; $\Delta$Saldo in der ersten Periode, Änderung des
Gini-Koeffizienten je 10~Mrd.~\euro{} Saldoverbesserung, Änderung des BIP je
Mrd.~\euro{} Saldoverbesserung in der ersten Periode und des BIP-Niveaus 2041
je Mrd.~\euro{} Saldoverbesserung in der letzten. Geordnet vom progressivsten zum
regressivsten Instrument.}
\label{tab:preise}
% Erzeugt von docs/modell/messung.js — nicht von Hand bearbeiten
\begin{tabular}{@{}lrrrr@{}}
\toprule
 & $\Delta$Saldo & Gini-Pkt. & \multicolumn{2}{c@{}}{BIP je Mrd.\,\euro{}} \\
\cmidrule(l){4-5}
Stellgröße (Schritt) & Mrd.\,\euro{} & je 10 Mrd.\,\euro{} & 2025--28 & Niveau 2041 \\
\midrule
Spitzensteuersatz (+1\,Pp.) & 0,1 & $-$1,65 & $-$0,85 & $-$0,63 \\
Mindeststeuer Milliardäre (+1\,Pp.) & 3,4 & $-$0,42 & $-$0,21 & 0,00 \\
Beitragsbemessungsgrenze (+10.000\,\euro{}) & 6,2 & $-$0,37 & $-$0,49 & 0,00 \\
Abgeltungsteuer (+1\,Pp.) & 2,3 & $-$0,32 & $-$0,36 & 0,00 \\
Erbschaftsteuer (+1\,Pp.) & 0,3 & $-$0,28 & $-$0,51 & 0,00 \\
Vermögensteuer (+0,1\,Pp.) & 2,8 & $-$0,27 & $-$0,49 & 0,00 \\
Bodenwertsteuer (+0,1\,Pp.) & 4,2 & $-$0,26 & $-$0,47 & 0,00 \\
Gewerbesteuer (+1\,Pp.) & 3,9 & $-$0,18 & $-$0,52 & $-$1,29 \\
Körperschaftsteuer (+1\,Pp.) & 2,0 & $-$0,17 & $-$0,51 & $-$4,78 \\
Eingangssteuersatz (+1\,Pp.) & 2,6 & $-$0,14 & $-$1,06 & $-$1,33 \\
Grundfreibetrag ($-$1.000\,\euro{}) & 8,1 & $+$0,05 & $-$0,93 & $-$0,36 \\
RV-Beitrag (+1\,Pp.) & 11,7 & $+$0,05 & $-$0,88 & 0,00 \\
Umsatzsteuer, ermäßigt (+1\,Pp.) & 4,4 & $+$0,07 & $-$0,67 & 0,00 \\
Umsatzsteuer (+1\,Pp.) & 7,7 & $+$0,07 & $-$0,72 & 0,00 \\
CO\textsubscript{2}-Preis (+10\,\euro{}/t) & 2,0 & $+$0,09 & $-$0,70 & 0,00 \\
KV-Beitrag (+1\,Pp.) & 12,1 & $+$0,10 & $-$0,79 & 0,00 \\
Rentenniveau ($-$1\,Pp.) & 5,8 & $+$0,38 & $-$0,83 & 0,00 \\
Kindergeld ($-$10\,\euro{}/Monat) & 1,7 & $+$0,40 & $-$0,80 & 0,00 \\
Bürgergeld ($-$10\,\euro{}/Monat) & 0,4 & $+$1,18 & $-$1,05 & 0,00 \\
\bottomrule
\end{tabular}

\end{table}

\subsection{Mechanische Wirkung, Verhalten und Nachfrage}
\label{sec:zerlegung}

Jede Saldowirkung der ersten Periode zerfällt genau in drei Teile, wenn man das
Modell stufenweise einschaltet:
\begin{equation}
\Delta S=\underbrace{\Delta S\big|_{\varepsilon=0,\;u=0}}_{\text{mechanisch}}
+\underbrace{\Bigl(\Delta S\big|_{u=0}-\Delta S\big|_{\varepsilon=0,\;u=0}\Bigr)}_{\text{Verhalten}}
+\underbrace{\Bigl(\Delta S-\Delta S\big|_{u=0}\Bigr)}_{\text{Nachfrage}} .
\end{equation}
Mechanisch heißt: alle Elastizitäten null und keine Nachfrage, also die Rechnung
mit festen Bemessungsgrundlagen. Verhalten ist, was Arbeitsangebot, Konsum,
Investitionen, Emissionen und Vermögen daran ändern, Nachfrage, was die veränderte
Kaufkraft der Haushalte über das BIP an Einnahmen kostet oder bringt. Die
Reihenfolge ist festgelegt; eine andere verschiebt kleine Beträge zwischen
Verhalten und Nachfrage, weil beide über das Nettoeinkommen verbunden sind
(Tabelle~\ref{tab:zerlegung}).

Beitragssätze verlieren rund ein Viertel ihrer mechanischen Wirkung über die
Nachfrage. Ein Prozentpunkt Rentenbeitrag bringt mechanisch \Wb{zerl}{rv}{mechanisch}~Mrd.~\euro,
die Nachfrage nimmt davon \Wb{zerl}{rv}{nachfrage}~Mrd.~\euro{} wieder weg, weil
Beiträge die Mitte und die unteren Einkommen treffen, deren Konsumneigung hoch
ist. Verhalten spielt keine Rolle, denn Beiträge gehen im Modell nicht in den
Grenzsteuersatz des Arbeitsangebots ein. Bei der Umsatzsteuer wirken beide
Kanäle: Von \Wb{zerl}{mwst}{mechanisch}~Mrd.~\euro{} mechanisch kostet die
Konsumreaktion \Wb{zerl}{mwst}{verhalten}, die Nachfrage
\Wb{zerl}{mwst}{nachfrage}~Mrd.~\euro.

Entlastungen finanzieren sich teilweise selbst. Ein um 1.000\,\euro{} höherer
Grundfreibetrag kostet mechanisch \Wb{zerl}{freibetrag}{mechanisch}~Mrd.~\euro;
mehr Arbeitsangebot bringt \Wb{zerl}{freibetrag}{verhalten}, die Nachfrage
\Wb{zerl}{freibetrag}{nachfrage}~Mrd.~\euro{} zurück, netto bleiben
\Wb{zerl}{freibetrag}{gesamt}~Mrd.~\euro. Den größten Teil der Angebotsreaktion
liefert D3: Sein mittleres zu versteuerndes Einkommen liegt knapp über dem
Grundfreibetrag und fällt bei 1.000\,\euro{} mehr ganz aus dem Tarif, sodass sein
Grenzsteuersatz auf null sinkt. D3 gewinnt deshalb
\Wb{dezil}{freibetrag}{D3}\,\euro, D4 nur \Wb{dezil}{freibetrag}{D4}\,\euro. Ein
Klimageld kostet mechanisch \Wb{zerl}{klimageld}{mechanisch}~Mrd.~\euro; über die
Nachfrage fließen \Wb{zerl}{klimageld}{nachfrage}~Mrd.~\euro{} zurück.

Verhalten zehrt dort, wo die Bemessungsgrundlage ausweichen kann. Bei der
Körperschaftsteuer kostet die Investitionsreaktion
\Wb{zerl}{kst}{verhalten}~Mrd.~\euro{} von \Wb{zerl}{kst}{mechanisch}~Mrd.~\euro{}
mechanisch; der größere Preis zeigt sich erst im Pfad (Abschnitt~\ref{sec:pfad}).
Der Spitzensteuersatz bringt mechanisch nur \Wb{zerl}{spitze}{mechanisch}~Mrd.~\euro,
weil er allein den Teil der Einkommen im obersten Prozent trifft, der im
Pareto-Rand über der Grenze liegt; das Ausweichen nimmt davon
\Wb{zerl}{spitze}{verhalten}~Mrd.~\euro{} zurück. Der \COz-Preis untergräbt
seine eigene Basis, wie beabsichtigt: Die Emissionsminderung kostet
\Wb{zerl}{co2}{verhalten}~Mrd.~\euro{} der mechanischen
\Wb{zerl}{co2}{mechanisch}~Mrd.~\euro.

\begin{table}[tbp]
\centering\footnotesize
\setlength{\tabcolsep}{4pt}
\caption{Zerlegung der Saldowirkung in der ersten Periode (2025--2028,
Jahreswert). Anteile in Prozent der Gesamtwirkung, nur wo diese mindestens
0,05~Mrd.~\euro{} beträgt. Die feste Ausweichquote der Mindeststeuer und die
Klammern der Reaktionen zählen zur mechanischen Wirkung.}
\label{tab:zerlegung}
% Erzeugt von docs/modell/messung.js — nicht von Hand bearbeiten
\begin{tabular}{lrrrrrrr}
\toprule
 & \multicolumn{4}{c}{Saldowirkung, Mrd.\,\euro{}/a} & \multicolumn{3}{c}{Anteil in \%} \\
\cmidrule(lr){2-5}\cmidrule(lr){6-8}
Stellgröße (Schritt) & mechan. & Verhalten & Nachfrage & gesamt & mech. & Verh. & Nachfr. \\
\midrule
Grundfreibetrag (1.000 \euro{}) & $-$11,05 & $+$1,04 & $+$3,16 & $-$6,86 & 161 & $-$15 & $-$46 \\
Eingangssteuersatz (1 Pp.) & $+$4,17 & $-$0,53 & $-$1,05 & $+$2,59 & 161 & $-$20 & $-$40 \\
Spitzensteuersatz (1 Pp.) & $+$0,38 & $-$0,06 & $-$0,03 & $+$0,30 & 130 & $-$19 & $-$11 \\
Umsatzsteuer (1 Pp.) & $+$10,61 & $-$0,88 & $-$2,06 & $+$7,67 & 138 & $-$12 & $-$27 \\
Erbschaftsteuer (1 Pp.) & $+$0,30 & 0,00 & $-$0,05 & $+$0,25 & 119 & 0 & $-$19 \\
Körperschaftsteuer (1 Pp.) & $+$2,88 & $-$0,45 & $-$0,39 & $+$2,04 & 141 & $-$22 & $-$19 \\
Gewerbesteuer (1 Pp.) & $+$5,09 & $-$0,45 & $-$0,75 & $+$3,89 & 131 & $-$12 & $-$19 \\
RV-Beitrag (1 Pp.) & $+$15,59 & 0,00 & $-$3,87 & $+$11,72 & 133 & 0 & $-$33 \\
KV-Beitrag (1 Pp.) & $+$15,72 & 0,00 & $-$3,58 & $+$12,14 & 130 & 0 & $-$30 \\
Rentenniveau (1 Pp.) & $-$7,54 & 0,00 & $+$1,77 & $-$5,77 & 131 & 0 & $-$31 \\
Bürgergeld (10 \euro{}/Monat) & $-$0,49 & 0,00 & $+$0,14 & $-$0,35 & 139 & 0 & $-$39 \\
Kindergeld (10 \euro{}/Monat) & $-$2,14 & 0,00 & $+$0,49 & $-$1,65 & 130 & 0 & $-$30 \\
CO\textsubscript{2}-Preis (10 \euro{}/t) & $+$3,20 & $-$0,62 & $-$0,53 & $+$2,05 & 157 & $-$31 & $-$26 \\
Klimageld (ein) & $-$12,59 & 0,00 & $+$2,88 & $-$9,71 & 130 & 0 & $-$30 \\
\midrule
Grenze Spitzensatz ($-$10.000 \euro{}) & $+$0,02 & $-$0,01 & 0,00 & $+$0,01 & -- & -- & -- \\
USt ermäßigt (1 Pp.) & $+$5,62 & $-$0,17 & $-$1,09 & $+$4,36 & 129 & $-$4 & $-$25 \\
Abgeltungsteuer (1 Pp.) & $+$2,58 & 0,00 & $-$0,31 & $+$2,27 & 114 & 0 & $-$14 \\
Beitragsbemessungsgr. (10.000 \euro{}) & $+$7,36 & 0,00 & $-$1,13 & $+$6,22 & 118 & 0 & $-$18 \\
Vermögensteuer (0,1 Pp.) & $+$3,33 & $-$0,07 & $-$0,50 & $+$2,75 & 121 & $-$3 & $-$18 \\
Bodenwertsteuer (0,1 Pp.) & $+$4,90 & 0,00 & $-$0,73 & $+$4,17 & 118 & 0 & $-$18 \\
Zucman-Mindeststeuer (1 Pp.) & $+$3,69 & 0,00 & $-$0,27 & $+$3,42 & 108 & 0 & $-$8 \\
\bottomrule
\end{tabular}

\end{table}

\subsection{Übertragungskanäle}
\label{sec:kanaele}

Abbildung~\ref{fig:struktur} zeigt, wie die Kanäle auf die Kennzahlen wirken,
unabhängig davon, welche Stellgröße sie bewegt hat. Ein Euro mehr Rente
verschlechtert den Saldo um einen Euro und erhöht das Nettoeinkommen eines
Haushalts um einen Euro. Mit Konsumneigung und Multiplikator werden daraus je
nach Dezil \WnachfrDzehnc{} bis \WnachfrDeins\,\euro{} Nachfrage, und jeder Euro
Nachfrage bringt \Wstabil\,\euro{} Einnahmen zurück. Ein Euro Rente an einen
Haushalt aus D1 kostet den Staat also nicht einen Euro, sondern
\WkostenDeins\,\euro. Über die Periode hinaus wirken nur Arbeitsangebot, Kapital
und Schulden.

Abbildung~\ref{fig:uebersicht} zeigt die gemessenen Kanten aller Stellgrößen am
Tisch. Jede Steuer bewegt mehr als ihren eigenen Kanal. Ein Prozentpunkt
Rentenbeitrag bringt \Wb{kanal}{rv}{sv}~Mrd.~\euro{} Beiträge, senkt aber das
Aufkommen aus Umsatzsteuer um \Wb{kanal}{rv}{mwst} und aus Einkommensteuer um
\Wb{kanal}{rv}{est}~Mrd.~\euro: die Umsatzsteuer, weil der Konsum mit dem
Nettoeinkommen sinkt und die Nachfrage nachlässt, die Einkommensteuer allein über
die Nachfrage. Der Nachfrageimpuls ist fast überall die zweitgrößte Kante; er ist
der Grund, warum keine Konsolidierung umsonst ist. Das Arbeitsangebot bewegen
messbar nur Grundfreibetrag und Eingangssatz, weil sie den Grenzsteuersatz der
Mitte ändern. Die vollständigen Messwerte stehen in
Tabelle~\ref{tab:kanaele}, die Ressortgraphen mit allen Kantengewichten in
Anhang~\ref{sec:anhang-graphen}.

\begin{figure}[tbp]
\centering
\resizebox{\textwidth}{!}{%
\begin{tikzpicture}[
  kn/.style={draw=grau, rounded corners=2pt, fill=hell, align=center, font=\small, text width=27mm, minimum height=10mm},
  zk/.style={draw=black, thick, rounded corners=2pt, fill=white, align=center, font=\small, text width=27mm, minimum height=10mm},
  pf/.style={-{Stealth[length=2.2mm]}, thick},
  rk/.style={-{Stealth[length=2.2mm]}, thick, dashed, grau},
  gl/.style={font=\scriptsize, fill=white, inner sep=1pt, align=center}]
% Spalte A: Kanäle
\node[kn] (aus)  at (0, 1.9)   {Ausgaben an Haushalte\\{\scriptsize Transfers, Renten}};
\node[kn] (ein)  at (0,-0.3)   {Einnahmen\\{\scriptsize ESt, USt, KSt, GewSt,\\Verm., \COz, Beiträge}};
\node[kn] (arb)  at (0,-5.8)   {Arbeitsangebot $\bar\lambda$};
\node[kn] (inv)  at (0,-7.7)   {Investitionsfaktor $I$};
\node[kn] (emi)  at (0,-9.8)   {Emissionen $E$};
% Spalte B: Haushalte und Nachfrage
\node[kn] (hh)   at (4.3,-2.4) {Nettoeinkommen\\der Haushalte $Y_i$};
\node[kn] (nf)   at (4.3,-4.7) {Nachfrageimpuls $\Delta Z$};
% Spalte C: Kennzahlen der Periode und Niveau
\node[zk] (sal)  at (8.6, 0.8) {Saldo $S$};
\node[zk] (ver)  at (8.6,-2.4) {Gini, Palma,\\Armutsrisiko};
\node[zk] (bip)  at (8.6,-4.7) {BIP der Periode};
\node[zk] (bipl) at (8.6,-6.9) {BIP-Niveau\\Folgeperiode};
\node[zk] (cob)  at (8.6,-9.8) {\COz-Budget\\(kumuliert)};
% Spalte D: Schulden
\node[zk] (sq)   at (12.9,-2.4) {Schuldenquote\\Folgeperiode};
% Kanten der Periode
\draw[pf] (ein.east) -- node[gl, pos=0.5]{$+1$} (sal.west);
\draw[pf] (aus.east) -- node[gl, pos=0.75]{$-1$} (sal.west);
\draw[pf] (ein.east) -- node[gl, pos=0.5, below left=0pt and 1pt]{$-1$ {\tiny Inzidenz}} (hh.west);
\draw[pf] (aus.east) -- node[gl, pos=0.18]{$+1$} (hh.west);
\draw[pf] (hh.east)  -- (ver.west);
\draw[pf] (hh.south) -- node[gl]{$\frac{0{,}6}{0{,}48}\,m_i$ {\tiny(\WnachfrDzehnc{} bis \WnachfrDeins)}} (nf.north);
\draw[pf] (nf.east)  -- node[gl]{$+1$} (bip.west);
\draw[pf] (nf.north east) -- node[gl, pos=0.78]{$+$\Wstabil} (sal.south west);
% Kanten über die Periode hinaus
\draw[pf] (sal.east) -- node[gl]{\Wschuld\,Pp.\\{\tiny je Mrd./a}} (sq.north west);
\draw[pf] (bipl.east) -- node[gl, pos=0.5]{Nenner} (sq.south);
\draw[pf] (arb.east) -- node[gl, pos=0.5]{$+$\Warbeit{} {\tiny Mrd. je \%}} (bipl.west);
\draw[pf] (inv.east) -- node[gl, pos=0.5]{$+$\Winvestn{} / \Winvestlang{} {\tiny Mrd. je \%, 1 Per. / lang}} (bipl.south west);
\draw[pf] (emi.east) -- node[gl]{$\frac{E_t}{\bar E(j_t)}\sum_{j\in t}\bar E(j)$ {\tiny entlang des Pfads}} (cob.west);
\draw[pf] (arb.north east) -- node[gl, pos=0.55]{$\tilde A_i=\lambda_iA_i$} (hh.south west);
\draw[rk] (sq.north) |- node[gl, pos=0.3, text=grau]{{\tiny Zinsen $z\,D_t$}\\{\tiny nächste Periode}} ($(sal.east)+(0,0.3)$);
\node[font=\scriptsize, text=grau, align=center] (ynote) at (8.6,-8.25) {$\times\,\mathit{BIP}_{t+1}/\mathit{BIP}_0$: alle Einnahmen\\der nächsten Periode};
\draw[rk] (bipl.south) -- (ynote.north);
% Rückwirkungen (gestrichelt)
\draw[rk] (bipl.north) -- node[gl, text=grau]{{\tiny Ausgangsniveau}\\{\tiny nächste Periode}} (bip.south);
\draw[rk] (arb.west) to[out=160,in=200,looseness=1.0] node[gl, pos=0.5, text=grau]{Basis} (ein.west);
\draw[rk] (inv.west) to[out=170,in=190,looseness=1.3] node[gl, pos=0.5, text=grau]{Basis $\bar I$} (ein.south west);
\draw[rk] (emi.west) to[out=175,in=185,looseness=1.5] node[gl, pos=0.5, text=grau]{\COz-Basis} ($(ein.south west)+(0,0.2)$);
\end{tikzpicture}}
\caption{Strukturgraph: wie Kanäle auf Kennzahlen wirken. Durchgezogene Kanten
tragen das Gewicht aus den Gleichungen des Abschnitts~\ref{sec:modell} im Status
quo, \euro{} Kennzahl je \euro{} Kanal, wenn nicht anders angegeben. Gestrichelte
Kanten sind Rückwirkungen auf die Bemessungsgrundlagen und in die nächste
Periode. $m_i$ ist die Grenzkonsumneigung, 0,15 in D10c bis 0,65 in D1.}
\label{fig:struktur}
\end{figure}

\begin{figure}[tbp]
\centering
\begin{tikzpicture}[
  x=1cm, y=1cm, font=\footnotesize,
  stell/.style={anchor=east, inner sep=2pt},
  kanal/.style={draw=grau, fill=hell, rounded corners=1.5pt, minimum width=27mm,
                minimum height=5mm, inner sep=1.5pt},
  ziel/.style={draw=grau, fill=white, rounded corners=1.5pt, minimum width=14mm,
               minimum height=6mm, font=\footnotesize\bfseries},
  struktur/.style={-{Stealth[length=1.6mm]}, grau},
  gewicht/.style={font=\scriptsize, fill=white, inner sep=0.8pt},
]
% Stellgrößen am Tisch, nach Ressort gruppiert (Schlüssel wie in messung.js)
\foreach \k/\t/\y in {freibetrag/Grundfreibetrag/0, eingang/Eingangssteuersatz/-0.6,
    spitze/Spitzensteuersatz/-1.2, mwst/Umsatzsteuer/-1.8, erb/Erbschaftsteuer/-2.4,
    kst/Körperschaftsteuer/-3.3, gewst/Gewerbesteuer/-3.9,
    rv/RV-Beitrag/-4.8, kv/KV-Beitrag/-5.4, rentenniveau/Rentenniveau/-6.0,
    bg/Bürgergeld/-6.6, kg/Kindergeld/-7.2, co2/\COz-Preis/-8.1, klimageld/Klimageld/-8.7}
  \node[stell] (s-\k) at (0,\y) {\t};
\foreach \r/\a/\b in {Finanzen/0/-2.4, Wirtschaft/-3.3/-3.9, Soziales/-4.8/-7.2, Umwelt/-8.1/-8.7}
  \draw[grau, thick] (-3.3,\a+0.2) -- (-3.3,\b-0.2)
    node[midway, left=1pt, rotate=90, anchor=south, font=\scriptsize\itshape] {\r};
% Kanäle (Spalten der Tabelle der Übertragungskanäle)
\foreach \k/\t/\y in {est/Einkommensteuer/0.2, arbeit/Arbeitsangebot/-0.75,
    mwst/Umsatzsteuer/-1.7, verm/Vermögensteuern/-2.65, unt/Unternehmensteuern/-3.6,
    invest/Investitionen/-4.55, nachfr/Nachfrage/-5.5, sv/Beiträge/-6.45,
    rente/Renten/-7.4, transfer/Transfers/-8.35, co2/\COz-Aufkommen/-9.3}
  \node[kanal] (k-\k) at (5.6,\y) {\t};
\node[ziel] (bip) at (10.3,-2.4) {BIP};
\node[ziel] (saldo) at (10.3,-6.0) {Saldo};
% Gemessene Kanten (generiert/uebersicht.tex)
\begin{scope}[on background layer]
% Erzeugt von docs/modell/messung.js — nicht von Hand bearbeiten
\ukante{freibetrag}{est}{neg}{2.15}
\ukante{freibetrag}{mwst}{pos}{0.87}
\ukante{freibetrag}{sv}{pos}{0.60}
\ukante{freibetrag}{arbeit}{pos}{2.32}
\ukante{freibetrag}{nachfr}{pos}{1.58}
\ukante{eingang}{est}{pos}{1.03}
\ukante{eingang}{mwst}{neg}{0.56}
\ukante{eingang}{sv}{neg}{0.46}
\ukante{eingang}{arbeit}{neg}{0.86}
\ukante{eingang}{nachfr}{neg}{0.79}
\ukante{mwst}{est}{neg}{0.46}
\ukante{mwst}{mwst}{pos}{1.71}
\ukante{mwst}{sv}{neg}{0.53}
\ukante{mwst}{nachfr}{neg}{1.16}
\ukante{erb}{verm}{pos}{0.44}
\ukante{kst}{unt}{pos}{0.75}
\ukante{kst}{invest}{neg}{2.60}
\ukante{kst}{nachfr}{neg}{0.54}
\ukante{gewst}{unt}{pos}{1.07}
\ukante{gewst}{sv}{neg}{0.45}
\ukante{gewst}{invest}{neg}{2.60}
\ukante{gewst}{nachfr}{neg}{0.68}
\ukante{rv}{est}{neg}{0.52}
\ukante{rv}{mwst}{neg}{0.72}
\ukante{rv}{sv}{pos}{2.58}
\ukante{rv}{nachfr}{neg}{1.83}
\ukante{kv}{est}{neg}{0.51}
\ukante{kv}{mwst}{neg}{0.70}
\ukante{kv}{sv}{pos}{2.60}
\ukante{kv}{nachfr}{neg}{1.72}
\ukante{rentenniveau}{est}{pos}{0.45}
\ukante{rentenniveau}{mwst}{pos}{0.44}
\ukante{rentenniveau}{sv}{pos}{0.51}
\ukante{rentenniveau}{rente}{pos}{1.44}
\ukante{rentenniveau}{nachfr}{pos}{1.06}
\ukante{bg}{transfer}{pos}{0.46}
\ukante{bg}{nachfr}{pos}{0.45}
\ukante{kg}{transfer}{pos}{0.68}
\ukante{kg}{nachfr}{pos}{0.58}
\ukante{co2}{co2}{pos}{0.76}
\ukante{co2}{nachfr}{neg}{0.60}
\ukante{klimageld}{est}{pos}{0.49}
\ukante{klimageld}{mwst}{pos}{0.47}
\ukante{klimageld}{sv}{pos}{0.58}
\ukante{klimageld}{co2}{neg}{2.14}
\ukante{klimageld}{transfer}{pos}{2.14}
\ukante{klimageld}{nachfr}{pos}{1.47}

\end{scope}
% Strukturelle Gewichte aus den Gleichungen (generiert/makros.tex)
\foreach \k in {est, mwst, verm, unt, sv, co2}
  \draw[struktur] (k-\k.east) to[out=0, in=180] (saldo.west);
\foreach \k in {rente, transfer}
  \draw[struktur, dashed] (k-\k.east) to[out=0, in=180] (saldo.west);
\draw[struktur] (k-nachfr.east) to[out=0, in=180] node[gewicht, pos=0.13] {\Wstabil} (saldo.west);
\draw[struktur] (k-nachfr.east) to[out=0, in=180] node[gewicht, pos=0.3] {1} (bip.west);
\draw[struktur] (k-arbeit.east) to[out=0, in=180] node[gewicht, pos=0.13] {\Warbeit} (bip.west);
\draw[struktur] (k-invest.east) to[out=0, in=180] node[gewicht, pos=0.13] {\Winvestn} (bip.west);
\end{tikzpicture}
\caption{Übertragungskanäle der Stellgrößen am Tisch in der ersten Periode.
Farbige Kanten sind gemessen (Tabelle~\ref{tab:kanaele}): blau, wenn der Kanal
steigt, orange, wenn er sinkt; die Strichstärke wächst mit dem Betrag, bezogen
auf die stärkste Wirkung derselben Einheit. Wirkungen unter 0,3~Mrd.~\euro{}
beziehungsweise 0,02\,\% entfallen. Graue Kanten folgen aus den Gleichungen:
Einnahmen gehen eins zu eins in den Saldo, Ausgaben (gestrichelt) mit
umgekehrtem Vorzeichen. Das Klimageld erscheint zweimal, als Transfer an die
Haushalte und als Minderung der \COz-Einnahmen, als die der Staat es bucht; in
den Saldo geht es einmal ein. Eine Milliarde Nachfrage hebt den Saldo um
\Wstabil~Mrd.~\euro, ein Prozent mehr Arbeitsangebot das BIP der folgenden
Periode um \Warbeit~Mrd.~\euro, ein um ein Prozent höherer Investitionsfaktor
das BIP nach einer Periode um \Winvestn~Mrd.~\euro.}
\label{fig:uebersicht}
\end{figure}

\subsection{Verteilung über die Haushaltstypen}

Tabelle~\ref{tab:dezile-wirkung} zeigt, bei wem die Wirkungen ankommen.
Umsatzsteuer und Beitragssätze belasten jeden Haushaltstyp, in Euro oben mehr,
im Verhältnis zum Einkommen unten. Die Bemessungsgrenzen deckeln die Beiträge,
die der Krankenversicherung ab D9, die der Rentenversicherung ab D10a: Ein
Prozentpunkt Rentenbeitrag kostet das oberste Prozent \Wb{dezil}{rv}{D10c}\,\euro,
eine um 10.000\,\euro{} höhere Beitragsbemessungsgrenze
\Wb{dezil}{bbg}{D10c}\,\euro; weil die Grenze der Krankenversicherung mitwandert,
zahlt auch D9. Die \COz-Abgabe kostet je 10\,\euro/t D1 \Wb{dezil}{co2}{D1}\,\euro{}
und D10c \Wb{dezil}{co2}{D10c}\,\euro; im Verhältnis zum Einkommen ist die Last
unten ein Vielfaches der Last oben. Das Klimageld zahlt allen gleich viel und
stellt D1 bis D3 besser, als sie ohne \COz-Preis stünden
(Tabelle~\ref{tab:dezile}). Das Rentenniveau kommt vor allem bei D3 bis D5 an,
das Bürgergeld nur bei D1 bis D4, das Kindergeld am meisten bei D3. Unternehmen-
und Vermögensteuern landen über die Inzidenzregel zu mehr als der Hälfte beim
obersten Dezil, die Mindeststeuer allein bei D10c.

\begin{table}[tbp]
\centering\footnotesize
\setlength{\tabcolsep}{2.25pt}
\caption{Änderung des verfügbaren Einkommens je Haushalt in der ersten
Periode, \euro{} je Jahr.}
\label{tab:dezile-wirkung}
% Erzeugt von docs/modell/messung.js — nicht von Hand bearbeiten
\begin{tabular}{@{}l*{12}{r}@{}}
\toprule
Stellgröße (Schritt) & D1 & D2 & D3 & D4 & D5 & D6 & D7 & D8 & D9 & D10a & D10b & D10c \\
\midrule
Grundfreibetrag (+1.000\,\euro{}) & 0 & 0 & $+$633 & $+$340 & $+$357 & $+$380 & $+$408 & $+$446 & $+$512 & $+$608 & $+$626 & $+$637 \\
Eingangssteuersatz (+1\,Pp.) & 0 & 0 & $-$36 & $-$70 & $-$100 & $-$136 & $-$175 & $-$220 & $-$281 & $-$325 & $-$317 & $-$322 \\
Spitzensteuersatz (+1\,Pp.) & 0 & 0 & 0 & 0 & 0 & 0 & 0 & 0 & 0 & 0 & 0 & $-$1.609 \\
Umsatzsteuer (+1\,Pp.) & $-$87 & $-$130 & $-$161 & $-$182 & $-$204 & $-$229 & $-$260 & $-$298 & $-$361 & $-$397 & $-$618 & $-$1.555 \\
Erbschaftsteuer (+1\,Pp.) & 0 & 0 & 0 & $-$1 & $-$2 & $-$3 & $-$5 & $-$9 & $-$16 & $-$21 & $-$38 & $-$177 \\
Körperschaftsteuer (+1\,Pp.) & $-$8 & $-$12 & $-$17 & $-$21 & $-$27 & $-$32 & $-$41 & $-$54 & $-$78 & $-$108 & $-$276 & $-$1.616 \\
Gewerbesteuer (+1\,Pp.) & $-$16 & $-$24 & $-$33 & $-$40 & $-$52 & $-$62 & $-$79 & $-$104 & $-$152 & $-$209 & $-$534 & $-$3.128 \\
RV-Beitrag (+1\,Pp.) & $-$129 & $-$194 & $-$248 & $-$304 & $-$366 & $-$440 & $-$527 & $-$649 & $-$842 & $-$974 & $-$979 & $-$987 \\
KV-Beitrag (+1\,Pp.) & $-$129 & $-$194 & $-$248 & $-$304 & $-$366 & $-$440 & $-$527 & $-$649 & $-$665 & $-$670 & $-$674 & $-$679 \\
Rentenniveau (+1\,Pp.) & $+$122 & $+$202 & $+$236 & $+$243 & $+$238 & $+$218 & $+$182 & $+$152 & $+$133 & $+$109 & $+$115 & $+$117 \\
Bürgergeld (+10\,\euro{}/Monat) & $+$72 & $+$30 & $+$10 & $+$2 & 0 & 0 & 0 & 0 & 0 & 0 & 0 & 0 \\
Kindergeld (+10\,\euro{}/Monat) & $+$45 & $+$61 & $+$67 & $+$64 & $+$59 & $+$53 & $+$47 & $+$42 & $+$36 & $+$28 & $+$20 & $+$11 \\
CO\textsubscript{2}-Preis (+10\,\euro{}/t) & $-$24 & $-$35 & $-$42 & $-$49 & $-$55 & $-$62 & $-$70 & $-$78 & $-$90 & $-$97 & $-$143 & $-$303 \\
Klimageld (ein) & $+$307 & $+$307 & $+$307 & $+$307 & $+$307 & $+$307 & $+$307 & $+$307 & $+$307 & $+$307 & $+$307 & $+$307 \\
\addlinespace
\multicolumn{13}{@{}l}{\emph{Weitere Stellgrößen}} \\
Grenze Spitzensatz ($-$10.000\,\euro{}) & 0 & 0 & 0 & 0 & 0 & 0 & 0 & 0 & 0 & 0 & 0 & $-$72 \\
Umsatzsteuer, ermäßigt (+1\,Pp.) & $-$46 & $-$69 & $-$85 & $-$96 & $-$108 & $-$122 & $-$138 & $-$158 & $-$191 & $-$210 & $-$327 & $-$824 \\
Abgeltungsteuer (+1\,Pp.) & $-$1 & $-$2 & $-$5 & $-$6 & $-$11 & $-$13 & $-$21 & $-$32 & $-$60 & $-$92 & $-$369 & $-$2.984 \\
Beitragsbemessungsgrenze (+10.000\,\euro{}) & 0 & 0 & 0 & 0 & 0 & 0 & 0 & 0 & $-$549 & $-$1.384 & $-$1.403 & $-$1.431 \\
Vermögensteuer (+0,1\,Pp.) & 0 & $-$1 & $-$4 & $-$9 & $-$19 & $-$32 & $-$53 & $-$91 & $-$165 & $-$227 & $-$400 & $-$1.868 \\
Bodenwertsteuer (+0,1\,Pp.) & 0 & $-$2 & $-$6 & $-$14 & $-$27 & $-$47 & $-$78 & $-$132 & $-$241 & $-$331 & $-$584 & $-$2.726 \\
Mindeststeuer Milliardäre (+1\,Pp.) & 0 & 0 & 0 & 0 & 0 & 0 & 0 & 0 & 0 & 0 & 0 & $-$9.476 \\
\bottomrule
\end{tabular}

\end{table}

\subsection{Langfristige Wirkungen}
\label{sec:pfad}

Am Ende des Pfads summieren sich die Saldowirkungen in der Schuldenquote
(Tabelle~\ref{tab:pfad} im Anhang). Wirkt eine Stellgröße nur über den Saldo,
senkt eine Milliarde Saldoverbesserung in der ersten Periode die Schuldenquote
2041 um rund 0,3~Prozentpunkte: Ein Prozentpunkt Umsatzsteuer senkt sie um
\Wb{pfad}{mwst}{schuld}, ein Prozentpunkt Rentenbeitrag um
\Wb{pfad}{rv}{schuld}~Punkte. Beim \COz-Preis und beim Klimageld ist es weniger,
weil das Aufkommen mit den Emissionen schrumpft. Bei Grundfreibetrag,
Eingangssatz und Unternehmensteuern bewegt ein zweiter Kanal das BIP im Nenner.
Der Grundfreibetrag kostet \Wb{kennzahl}{freibetrag}{saldo}~Mrd.~\euro{}
im Jahr, erhöht die Quote aber nur um \Wb{pfad}{freibetrag}{schuld}~Punkte, weil
das zusätzliche Arbeitsangebot das BIP-Niveau 2041 um
\Wb{pfad}{freibetrag}{bip}~Mrd.~\euro{} hebt; den größten Teil davon liefert D3
an der Tarifgrenze (Abschnitt~\ref{sec:zerlegung}). Die Körperschaftsteuer bringt
\Wb{kennzahl}{kst}{saldo}~Mrd.~\euro, senkt die Quote aber nur um
\Wb{pfad}{kst}{schuld}~Punkte, weil das private Kapital schrumpft, das sich mit
7\,\% je Jahr anpasst. Der \COz-Preis ist die einzige Stellgröße am Tisch, die die
kumulierten Emissionen senkt, um \Wb{pfad}{co2}{co2kum}~Mt je 10\,\euro/t bis
2040.

\subsection{Nichtlinearität}
\label{sec:nichtlinear}

Abbildung~\ref{fig:kurven} verschiebt fünf Stellgrößen über ihren ganzen
Spielraum. Eine Gerade heißt, dass jeder Schritt gleich wirkt, ein Knick, dass
eine Schwelle oder Klammer greift, eine flacher werdende Kurve, dass die
Verhaltensreaktion einen wachsenden Teil der mechanischen Wirkung aufzehrt.
Unterhalb von 42\,\% deckelt der Spitzensatz auch die Proportionalzone, und eine
Senkung wird zur breiten Steuersenkung; darüber trifft er nur den Pareto-Rand, und
ab 45 und 60\,\% flachen Ausweichen und Wegzug die Kurve ab. Ihr Maximum liegt
oberhalb des Spielraums. Die Formel von \textcite{saez2001} für den
aufkommensmaximierenden Spitzensatz, $1/(1+\alpha e)$, ergäbe mit $\alpha = 1{,}5$
und einer Elastizität des zu versteuernden Einkommens von $e = 0{,}25$
\parencite{saez2012} 73\,\%. Im Modell ist das Ausweichen linear im Satz, eine
Semi-Elastizität; die Elastizität sinkt deshalb mit steigendem Satz, und das
Maximum liegt höher. Die Umsatzsteuer bringt je Punkt umso weniger, je höher sie
steht, weil der Anteil $\tau/(1+\tau)$ konkav ist und der Konsum reagiert. Die
Körperschaftsteuer ist Satz mal schrumpfende Basis, eine nach unten geöffnete
Parabel mit dem Scheitel weit außerhalb des Spielraums. Das Rentenniveau wirkt
linear. Beim \COz-Preis fallen die Emissionen linear, solange die Reaktion
zwischen ihren Klammern liegt, und erreichen bei \COzBoden\,\euro/t die
Untergrenze. Aufkommen und Saldo erreichen bei \COzMax\,\euro/t ein Maximum und
fallen dann, weil die Bemessungsgrundlage schneller schrumpft, als der Preis
steigt; jenseits der Untergrenze steigen sie wieder linear mit dem Preis, ohne
weitere Minderung.

Auch um den Status quo sind nicht alle Wirkungen symmetrisch
(Tabelle~\ref{tab:symmetrie}). Der Spitzensatz steht genau an der Schwelle, ab der
das Ausweichen einsetzt: Ein Punkt mehr bringt \Wb{sym}{spitze}{auf}~Mrd.~\euro, ein
Punkt weniger kostet \Wb{sym}{spitze}{ab}. Ein höherer Grundfreibetrag kostet
\Wb{sym}{freibetrag}{auf}~Mrd.~\euro, ein niedrigerer bringt
\Wb{sym}{freibetrag}{ab}: Beim höheren fällt D3 aus dem Tarif, spart nur seine
geringe Steuer und arbeitet mehr, beim niedrigeren zahlt D3 auf die vollen
1.000\,\euro. Tabelle~\ref{tab:preise} misst deshalb jede Stellgröße in ihrer
Konsolidierungsrichtung.

\newcommand{\kurve}[6]{% Stellgröße, Spalte, x-Beschriftung, y-Beschriftung, Titel, Farbe
  \nextgroupplot[title={#5}, xlabel={#3}, ylabel={#4},
    extra x ticks={\csname Wsq@#1\endcsname}, extra x tick labels={},
    extra x tick style={grid=major, major grid style={dashed, grau}}]
  \addplot[color=#6, very thick, mark=none] table[x=x, y=#2] {generiert/verlauf_#1.dat};}

\begin{figure}[p]
\centering
\begin{tikzpicture}
\begin{groupplot}[group style={group size=2 by 5, horizontal sep=18mm, vertical sep=19mm},
  width=0.44\textwidth, height=0.15\textheight,
  tick label style={font=\scriptsize}, label style={font=\scriptsize},
  title style={font=\small}, grid=none, axis line style={grau},
  /pgf/number format/use comma, /pgf/number format/1000 sep={.},
  scaled y ticks=false]
\kurve{spitze}{saldo}{Spitzensteuersatz, \%}{Saldo, Mrd.\,\euro}{Spitzensteuersatz: Saldo}{pos}
\kurve{spitze}{gini}{Spitzensteuersatz, \%}{Gini, Pkt.}{Spitzensteuersatz: Gini}{neg}
\kurve{mwst}{saldo}{Umsatzsteuer, \%}{Saldo, Mrd.\,\euro}{Umsatzsteuer: Saldo}{pos}
\kurve{mwst}{gini}{Umsatzsteuer, \%}{Gini, Pkt.}{Umsatzsteuer: Gini}{neg}
\kurve{kst}{saldo}{Körperschaftsteuer, \%}{Saldo, Mrd.\,\euro}{Körperschaftsteuer: Saldo}{pos}
\kurve{kst}{bip}{Körperschaftsteuer, \%}{BIP, Mrd.\,\euro}{Körperschaftsteuer: BIP der Periode}{neg}
\kurve{rentenniveau}{saldo}{Rentenniveau, \%}{Saldo, Mrd.\,\euro}{Rentenniveau: Saldo}{pos}
\kurve{rentenniveau}{gini}{Rentenniveau, \%}{Gini, Pkt.}{Rentenniveau: Gini}{neg}
\kurve{co2}{saldo}{\COz-Preis, \euro/t}{Saldo, Mrd.\,\euro}{\COz-Preis: Saldo}{pos}
\kurve{co2}{emiss}{\COz-Preis, \euro/t}{Emissionen, Mt/a}{\COz-Preis: Emissionen}{neg}
\end{groupplot}
\end{tikzpicture}
\caption{Saldo und eine zweite Kennzahl der ersten Periode über den ganzen
Spielraum der Stellgröße, alle übrigen im Status quo. Gestrichelt: Wert des Status
quo. 21 Stützstellen je Kurve.}
\label{fig:kurven}
\end{figure}

\begin{table}[tbp]
\centering\footnotesize
\caption{Nichtlinearität: Saldowirkung einer Erhöhung um den Schritt und einer
Senkung um denselben Schritt (mit umgekehrtem Vorzeichen), erste Periode,
Mrd.~\euro. Gleiche Werte heißen: in diesem Bereich linear.}
\label{tab:symmetrie}
% Erzeugt von docs/modell/messung.js — nicht von Hand bearbeiten
\begin{tabular}{lrrr}
\toprule
Stellgröße (Schritt) & $+\Delta$ & $-\Delta$ & Abweichung \\
\midrule
Grundfreibetrag (1.000 \euro{}) & $-$6,86 & $-$8,10 & 18\,\% \\
Eingangssteuersatz (1 Pp.) & $+$2,59 & $+$2,60 & linear \\
Spitzensteuersatz (1 Pp.) & $+$0,30 & $+$0,30 & linear \\
Umsatzsteuer (1 Pp.) & $+$7,67 & $+$7,91 & 3\,\% \\
Erbschaftsteuer (1 Pp.) & $+$0,25 & $+$0,25 & linear \\
Körperschaftsteuer (1 Pp.) & $+$2,04 & $+$2,06 & linear \\
Gewerbesteuer (1 Pp.) & $+$3,89 & $+$3,93 & linear \\
RV-Beitrag (1 Pp.) & $+$11,72 & $+$11,79 & linear \\
KV-Beitrag (1 Pp.) & $+$12,14 & $+$12,21 & linear \\
Rentenniveau (1 Pp.) & $-$5,77 & $-$5,77 & linear \\
Bürgergeld (10 \euro{}/Monat) & $-$0,35 & $-$0,35 & linear \\
Kindergeld (10 \euro{}/Monat) & $-$1,65 & $-$1,65 & linear \\
CO\textsubscript{2}-Preis (10 \euro{}/t) & $+$2,05 & $+$2,20 & 7\,\% \\
Grenze Spitzensatz ($-$10.000 \euro{}) & $+$0,01 & $+$0,01 & linear \\
USt ermäßigt (1 Pp.) & $+$4,36 & $+$4,50 & 3\,\% \\
Abgeltungsteuer (1 Pp.) & $+$2,27 & $+$2,27 & linear \\
Beitragsbemessungsgr. (10.000 \euro{}) & $+$6,22 & $+$5,80 & 7\,\% \\
Vermögensteuer (0,1 Pp.) & $+$2,75 & -- & Status quo am Rand \\
Bodenwertsteuer (0,1 Pp.) & $+$4,17 & $+$4,17 & linear \\
Zucman-Mindeststeuer (1 Pp.) & $+$3,42 & -- & Status quo am Rand \\
\bottomrule
\end{tabular}

\end{table}
